% ============================================================================
% preamble.tex: оформление по указанията на гл. ас. д-р Боян Бележков за курсовия
% проект: Times New Roman 12 pt, разстояние 1,5 реда, бележки под линия с единично
% разстояние, двустранно подравнен текст, номерирани страници, български език.
% Двигател: pdfLaTeX с T2A и tempora (шрифт Times с пълна кирилица), както в
% шаблона octet/docs.
% ============================================================================

\usepackage{cmap}
\usepackage[T2A]{fontenc}
\usepackage[utf8]{inputenc}
\usepackage[bulgarian]{babel}
\usepackage{tempora}
\usepackage[a4paper, top=2cm, bottom=2.2cm, left=2cm, right=2cm, footskip=1cm]{geometry}

% 1,5 реда в основния текст. Пакетът setspace оставя бележките под линия
% с единично разстояние, както изискват указанията.
\usepackage{setspace}
\setstretch{1.5}
\usepackage[expansion=false]{microtype}
\usepackage{indentfirst}
\setlength{\parindent}{1.25cm}
\setlength{\parskip}{0pt}

% Бележки под линия: 10 pt, единично разстояние, малък отстъп между бележките.
\setlength{\footnotesep}{0.9\baselineskip}
\setlength{\skip\footins}{1.2em plus 0.3em minus 0.2em}

% Номера на страниците: долу в средата.
\usepackage{fancyhdr}
\pagestyle{fancy}
\fancyhf{}
\renewcommand{\headrulewidth}{0pt}
\fancyfoot[C]{\thepage}
\fancypagestyle{plain}{\fancyhf{}\renewcommand{\headrulewidth}{0pt}\fancyfoot[C]{\thepage}}

% Заглавия: арабска номерация, получер шрифт.
\usepackage{titlesec}
\titleformat{\section}{\fontsize{14}{17}\selectfont\bfseries}{\thesection.}{0.5em}{}
\titleformat{\subsection}{\fontsize{12}{15}\selectfont\bfseries}{\thesubsection.}{0.5em}{}
\titleformat{\subsubsection}{\fontsize{12}{15}\selectfont\bfseries\itshape}{\thesubsubsection.}{0.5em}{}
\titlespacing*{\section}{0pt}{1.1em}{0.5em}
\titlespacing*{\subsection}{0pt}{0.9em}{0.3em}
\titlespacing*{\subsubsection}{0pt}{0.7em}{0.2em}

\usepackage[titles]{tocloft}
\renewcommand{\cftsecleader}{\cftdotfill{\cftdotsep}}
\renewcommand{\cftsecaftersnum}{.}
\renewcommand{\cftsubsecaftersnum}{.}
\renewcommand{\cftsubsubsecaftersnum}{.}

% Надписи: таблиците отгоре, фигурите отдолу.
\usepackage{caption}
\DeclareCaptionLabelFormat{tablelabel}{Таблица #2}
\captionsetup[table]{labelformat=tablelabel, labelsep=period, position=above, font={small}, labelfont={bf}, skip=4pt}
\DeclareCaptionLabelFormat{figurelabel}{Фигура #2}
\captionsetup[figure]{labelformat=figurelabel, labelsep=period, position=below, font={small}, labelfont={bf}, skip=4pt}
\newcommand{\tabref}[1]{(табл.~\ref{#1})}
\newcommand{\figref}[1]{(фиг.~\ref{#1})}

\usepackage{graphicx}
\usepackage{float}
\usepackage{array}
\usepackage{booktabs}
\usepackage{tabularx}
\usepackage{longtable}
\usepackage{ragged2e}
% \hspace{0pt} позволява пренасяне и на първата дума в клетката.
\newcolumntype{L}[1]{>{\RaggedRight\arraybackslash\hspace{0pt}}p{#1}}
\newcolumntype{Y}{>{\RaggedRight\arraybackslash\hspace{0pt}}X}
% Колона за термините в речника: получер шрифт, без пренасяне.
\newcolumntype{T}[1]{>{\RaggedRight\arraybackslash\bfseries\hyphenpenalty=10000\exhyphenpenalty=10000}p{#1}}
\usepackage[table]{xcolor}
\usepackage{tikz}
\usetikzlibrary{arrows.meta,positioning,shapes.geometric,fit,calc}
\usepackage{enumitem}
\setlist{nosep, leftmargin=1.6em}
\usepackage{multicol}
% Списък със съкращения в две колони: съкращението вляво, обяснението вдясно.
\newlist{abbrlist}{description}{1}
\setlist[abbrlist]{style=sameline, font=\normalfont, leftmargin=1.9cm, labelwidth=1.7cm,
  labelsep=0.2cm, itemsep=0.1em, parsep=0pt, topsep=0.2em}
\usepackage{textcomp}
\usepackage{xurl}
\usepackage[hidelinks]{hyperref}
\urlstyle{same}
\hypersetup{
  pdftitle={Ролята на прокурора в досъдебната фаза на наказателния процес},
  pdfauthor={Алекс Иванов Цветанов},
  pdfsubject={Курсов проект по Процесуално-правни аспекти на киберпрестъпленията},
  pdflang={bg-BG},
}

% Двустранно подравнен текст с дълги наименования на актове има нужда от място.
\setlength{\emergencystretch}{3em}

% Таблици и фигури с единично разстояние.
\usepackage{etoolbox}
\AtBeginEnvironment{tabular}{\setstretch{1.1}}
\AtBeginEnvironment{tabularx}{\setstretch{1.1}}
\AtBeginEnvironment{longtable}{\setstretch{1.1}}
\AtBeginEnvironment{tikzpicture}{\setstretch{1.0}}

\tikzset{
  box/.style={draw, rounded corners=2pt, align=center, font=\small, minimum height=0.9cm, inner sep=4pt, fill=blue!4},
  court/.style={box, fill=orange!10},
  done/.style={box, fill=green!7},
  arr/.style={-Latex, thick},
}

% Номерация на библиографията, която продължава между отделните групи източници.
\newlist{biblist}{enumerate}{1}
\setlist[biblist]{label={[\arabic*]}, leftmargin=2.6em, labelsep=0.5em, itemsep=0pt, align=left}
